\chapter{附加数学工具：Bessel 平均的一个矩展开命题}
\label{app:analytic-series}

正文的实验概率已经统一为
\begin{equation}
P_n
=\int\dd\lambda\,W(\lambda)
J_n^2\!\left(2|\beta(\lambda)|\right).
\label{eq:analytic-series-master}
\end{equation}
如果数值积分足够，本附录可以全部跳过。这里唯一要做的事情，是在局域耦合可以写成
\begin{equation}
\beta(\lambda)=\beta_0f(\lambda),
\qquad |f(\lambda)|\le1
\end{equation}
时，把同一个非线性平均改写成 \(|f|\) 的矩。

\section{一个命题：非线性 Bessel 平均等于矩序列的线性组合}

对 \(n\ge0\)，
\begin{equation}
J_n(2x)
=\sum_{j=0}^{\infty}
\frac{(-1)^j}{j!(n+j)!}x^{n+2j}.
\label{eq:Bessel-series-definition}
\end{equation}
定义纯数学系数
\begin{equation}
 a_j^{(n)}\equiv\frac{(-1)^j}{j!(n+j)!},
\end{equation}
则
\begin{align}
J_n^2\!\left(2|\beta_0||f|\right)
&=\left|\sum_{j\ge0}a_j^{(n)}|\beta_0|^{n+2j}|f|^{n+2j}\right|^2\\
&=\sum_{j,k\ge0}
a_j^{(n)}a_k^{(n)}
|\beta_0|^{2(n+j+k)}|f|^{2(n+j+k)}.
\end{align}
因此只要允许交换求和与积分，
\begin{equation}
\boxed{
P_n
=\sum_{j,k\ge0}
a_j^{(n)}a_k^{(n)}
|\beta_0|^{2S_{njk}}M_{S_{njk}},
\qquad
S_{njk}\equiv n+j+k,}
\label{eq:moment-expansion-master}
\end{equation}
其中
\begin{equation}
\boxed{
M_s\equiv\int\dd\lambda\,W(\lambda)|f(\lambda)|^{2s}.}
\label{eq:moment-definition}
\end{equation}
这就是原来所有双重级数的最简数学结构。旧记号中的 \(C_j^n\) 只是在 \(a_j^{(n)}\) 上再吸收 \(\beta_0\) 和固定相位；第二个指标 \(k\) 只来自取模平方；所谓 \(D_{jk}^n\) 如果只依赖 \(n+j+k\)，就只是同一列矩 \(M_s\) 的重新命名。

对有限 \(|\beta_0|\)，Bessel 级数绝对收敛；在 \(|f|\le1\) 且 \(W\) 为归一化有限测度时，可以用一致控制交换求和与积分。若 \(|f|\) 不有界，则必须单独检查各阶矩 \(M_s\) 是否存在。

\section{两个算例：Gaussian 时间矩与指数横向矩}

正文的 Gaussian 延迟问题对应
\begin{equation}
f_t(t;\tau)
=\exp\!\left[-\frac{(t-\tau)^2}{4\sigma_{\rm L}^2}\right],
\qquad
W_t(t)=\frac{1}{\sqrt{2\pi}\sigma_{\rm e}}
\exp\!\left[-\frac{t^2}{2\sigma_{\rm e}^2}\right].
\end{equation}
因此
\begin{align}
M_s^{(t)}(\tau)
&=\int_{-\infty}^{+\infty}
\frac{\dd t}{\sqrt{2\pi}\sigma_{\rm e}}
\exp\!\left[-\frac{t^2}{2\sigma_{\rm e}^2}
-\frac{s(t-\tau)^2}{2\sigma_{\rm L}^2}\right]\\
&=\boxed{
\frac{1}{\sqrt{1+s(\sigma_{\rm e}/\sigma_{\rm L})^2}}
\exp\!\left[
-\frac{s(\tau/\sigma_{\rm L})^2}
{2[1+s(\sigma_{\rm e}/\sigma_{\rm L})^2]}
\right].}
\label{eq:gaussian-time-moment}
\end{align}
把它直接代入式~\eqref{eq:moment-expansion-master}，就得到与正文 Bessel--Gaussian 积分~\eqref{eq:Pn-gaussian} 完全等价的双重级数。弱耦合固定 \(n>0\) 时只保留 \(j=k=0\)，即 \(s=n\)，于是立即恢复正文
\begin{equation}
\sigma_{\tau,n}^2
=\sigma_{\rm e}^2+\frac{\sigma_{\rm L}^2}{n}.
\end{equation}

空间平均完全同理。若
\begin{equation}
\beta(\bm R_\perp,t)
=\beta_0f_\perp(\bm R_\perp)f_t(t),
\end{equation}
则
\begin{equation}
M_s=M_s^{(\perp)}M_s^{(t)},
\qquad
M_s^{(\perp)}
=\frac{\displaystyle\int_{\mathcal A}\dd^2R_\perp\,
W_\perp|f_\perp|^{2s}}
{\displaystyle\int_{\mathcal A}\dd^2R_\perp\,W_\perp}.
\label{eq:general-spatial-moment}
\end{equation}
旧文中的空间系数最自然地就是
\begin{equation}
\boxed{D_s\equiv M_s^{(\perp)}.}
\label{eq:Dfactor_start}
\end{equation}

例如圆形区域 \(a\le b\le w\) 上取
\begin{equation}
f_\perp(b,\varphi)
=\ee^{-(b-a)/\delta}\cos\varphi,
\label{eq:F_spatial_model}
\end{equation}
并采用均匀面积权重，则
\begin{equation}
D_s
=\frac{1}{\pi(w^2-a^2)}
\int_0^{2\pi}\dd\varphi\,\cos^{2s}\varphi
\int_a^w b\dd b\,\ee^{-2s(b-a)/\delta}.
\end{equation}
其中
\begin{equation}
\int_0^{2\pi}\cos^{2s}\varphi\,\dd\varphi
=2\sqrt\pi\,\frac{\Gamma(s+1/2)}{\Gamma(s+1)},
\end{equation}
所以对 \(s>0\)
\begin{equation}
\boxed{
D_s
=\frac{\delta\left[(2sa+\delta)-
\ee^{-2s(w-a)/\delta}(2sw+\delta)\right]}
{2\sqrt\pi\,s^2(w^2-a^2)}
\frac{\Gamma(s+1/2)}{\Gamma(s+1)},
\qquad D_0=1.}
\label{eq:D_sphere}
\end{equation}
Gamma 函数只来自 \(\cos^{2s}\varphi\) 的角向积分。最终
\begin{equation}
P_{\rm EEGS}(n)
=\sum_{j,k\ge0}
a_j^{(n)}a_k^{(n)}
|\beta_0|^{2S_{njk}}
D_{S_{njk}}M_{S_{njk}}^{(t)},
\end{equation}
与正文的空间--时间积分完全等价。

因此这些级数没有增加新的 PINEM 模型。它们只是在需要解析闭式时，把
\(\langle J_n^2(2|\beta|)\rangle\)
改写成 \(|\beta|\) 的高阶矩；不需要闭式时，正文的直接积分更短，也更接近物理本体。
